\section{产品拆解：Lovart 应该像工具、平台还是工作台}

前面讲了定位、工作流和壁垒，本章把 Lovart 放进产品形态里拆解。AI 设计产品有三种常见形态：工具、平台、工作台。工具解决单点任务，比如生成海报或改图；平台连接创作者、素材和分发；工作台则承载一个专业角色的日常生产过程。Lovart 如果要成为垂直 Agent，更像工作台，而不是单点工具。

\begin{knowledgebox}{工具、平台、工作台}
\begin{tabularx}{\textwidth}{p{0.20\textwidth}p{0.34\textwidth}X}
\toprule
形态 & 用户为什么来 & 壁垒来源 \\
\midrule
工具 & 快速完成单点任务 & 功能体验和低成本。 \\
平台 & 找素材、找模板、找用户或分发 & 网络效应和内容供给。 \\
工作台 & 管理项目、资产、协作和交付 & 工作流沉淀和组织迁移成本。 \\
\bottomrule
\end{tabularx}
\end{knowledgebox}

\begin{importantbox}{Lovart 更适合做工作台}
如果 Lovart 只做工具，它会被模型公司快速功能化；如果做平台，它要面对供给和分发难题；如果做工作台，它可以围绕设计师的连续工作流形成更深壁垒。
\end{importantbox}

\subsection{工作台里的四个闭环}

设计工作台需要四个闭环：项目闭环、资产闭环、反馈闭环和交付闭环。项目闭环让 brief、任务和进度可追踪；资产闭环让品牌素材、参考图和生成结果可复用；反馈闭环让客户/团队意见进入下一轮迭代；交付闭环让最终文件可编辑、可交付、可商用。

\begin{knowledgebox}{设计工作台闭环表}
\begin{tabularx}{\textwidth}{p{0.20\textwidth}p{0.34\textwidth}X}
\toprule
闭环 & 要解决什么 & 没有它的后果 \\
\midrule
项目闭环 & brief、任务、进度、版本 & 用户只能做临时生成。 \\
资产闭环 & 品牌素材、历史方案、参考图 & 每次都从零开始，难以专业化。 \\
反馈闭环 & 修改意见和评审记录 & Agent 不理解团队偏好。 \\
交付闭环 & 文件格式、版权、尺寸、协作 & 作品不能进入真实生产。 \\
\bottomrule
\end{tabularx}
\end{knowledgebox}

\subsection{本章小结}

Lovart 的产品上限取决于它能否从工具走向工作台。垂直 Agent 的护城河不在一次生成，而在长期项目和资产闭环。

